% Liaisons polarisées et polarité de trois molécules (L'essentiel, chapitre 6)
\begin{tikzpicture}[x=1cm,y=1cm,
  atome/.style={font=\otSemi\fontsize{10}{11}\selectfont,text=otInk,inner sep=1.5pt},
  lien/.style={otInk!80,line width=0.9pt},
  bary/.style={font=\otMedium\fontsize{7}{8}\selectfont,inner sep=1pt},
  nomL/.style={pcLegende,anchor=north,align=center}]
% Chlorure d'hydrogène
\node[atome] (h) at (0,0) {H}; \node[atome] (cl) at (1.3,0) {Cl};
\draw[lien] (h) -- (cl);
\node[pcDelta] at (0,0.42) {\dplus}; \node[pcDelta] at (1.3,0.42) {\dmoins};
\node[nomL] at (0.65,-0.75) {$\chi(\ce{Cl})>\chi(\ce{H})$ :\\liaison polarisée,\\molécule \textbf{polaire}};
% Eau
\begin{scope}[shift={(4.6,0.3)}]
  \node[atome] (o) at (0,0) {O}; \node[atome] (h1) at (-0.95,-0.73) {H}; \node[atome] (h2) at (0.95,-0.73) {H};
  \draw[lien] (o) -- (h1) (o) -- (h2);
  \node[pcDelta] at (0,0.42) {\textcolor{otNavy}{$2\delta^{-}$}}; \node[pcDelta] at (-1.3,-0.4) {\dplus}; \node[pcDelta] at (1.3,-0.4) {\dplus};
  \draw[otRed!80!black,line width=0.6pt] (-0.07,-0.8) -- (0.07,-0.66) (-0.07,-0.66) -- (0.07,-0.8);
  \node[bary,text=otRed!80!black,anchor=north] at (0,-0.82) {G$^{+}$};
  \node[bary,text=otNavy,anchor=west] at (0.22,0.1) {G$^{-}$ (sur O)};
  \node[nomL] at (0,-1.05) {G$^{+}$ et G$^{-}$ distincts :\\molécule \textbf{polaire}};
\end{scope}
% Dioxyde de carbone
\begin{scope}[shift={(9.0,0)}]
  \node[atome] (o1) at (-1.15,0) {O}; \node[atome] (c) at (0,0) {C}; \node[atome] (o2) at (1.15,0) {O};
  \draw[lien] ([yshift=1.3pt]o1.east) -- ([yshift=1.3pt]c.west) ([yshift=-1.3pt]o1.east) -- ([yshift=-1.3pt]c.west)
    ([yshift=1.3pt]c.east) -- ([yshift=1.3pt]o2.west) ([yshift=-1.3pt]c.east) -- ([yshift=-1.3pt]o2.west);
  \node[pcDelta] at (-1.15,0.42) {\dmoins}; \node[pcDelta] at (0,0.42) {\textcolor{otRed!80!black}{$2\delta^{+}$}}; \node[pcDelta] at (1.15,0.42) {\dmoins};
  \node[bary,text=otInk,anchor=north] at (0,-0.22) {G$^{+}$ = G$^{-}$};
  \node[nomL] at (0,-0.75) {liaisons polarisées mais\\G$^{+}$ et G$^{-}$ confondus :\\molécule \textbf{apolaire}};
\end{scope}
\end{tikzpicture}
\par\smallskip\legende{G$^{+}$ et G$^{-}$ désignent les centres géométriques des charges partielles positives et négatives.}
